% Zoznam skratiek a značiek pre bioinformatiku a biológiu
% Tento súbor obsahuje základné skratky používané v bioinformatike a biológii

\newacronym{DNA}{DNA}{deoxyribonukleová kyselina}
\newacronym{RNA}{RNA}{ribonukleová kyselina}
\newacronym{PCR}{PCR}{polymerázová reťazová reakcia}
\newacronym{NGS}{NGS}{sekvenovanie novej generácie (Next Generation Sequencing)}
\newacronym{MSA}{MSA}{viacnásobné zarovnanie sekvencií (Multiple Sequence Alignment)}
\newacronym{SNP}{SNP}{jednonukleotidový polymorfizmus (Single Nucleotide Polymorphism)}
\newacronym{ORF}{ORF}{otvorený čítací rámec (Open Reading Frame)}
\newacronym{BLAST}{BLAST}{Basic Local Alignment Search Tool}
\newacronym{FASTA}{FASTA}{formát pre zápis sekvencií}
\newacronym{FASTQ}{FASTQ}{formát pre zápis sekvencií s kvalitou}
\newacronym{BAM}{BAM}{binárny formát pre zarovnané sekvencie}
\newacronym{SAM}{SAM}{sekvenčný formát pre zarovnané sekvencie}
\newacronym{GFF}{GFF}{General Feature Format}
\newacronym{CDS}{CDS}{kódujúca sekvencia (Coding Sequence)}
\newacronym{tRNA}{tRNA}{transferová RNA}
\newacronym{rRNA}{rRNA}{ribozomálna RNA}
\newacronym{mRNA}{mRNA}{messenger RNA}
\newacronym{bp}{bp}{bázový pár (base pair)}
\newacronym{kb}{kb}{kilobázový pár (kilobase pair)}
\newacronym{Mb}{Mb}{megabázový pár (megabase pair)}
\newacronym{aa}{aa}{aminokyselina}
\newacronym{PAM}{PAM}{protospacer adjacent motif}
\newacronym{CRISPR}{CRISPR}{Clustered Regularly Interspaced Short Palindromic Repeats}
\newacronym{Cas}{Cas}{CRISPR-associated protein}
\newacronym{GC}{GC}{obsah guanínu a cytozínu (GC content)}
\newacronym{UTR}{UTR}{neprekladaná oblasť (Untranslated Region)}
\newacronym{ID}{ID}{identifikátor}
